\documentclass[]{ceurart}

\sloppy

\usepackage{listings}
\lstset{breaklines=true}
\usepackage{booktabs}
\usepackage{amsmath}

\begin{document}

\copyrightyear{2026}
\copyrightclause{Copyright for this paper by its authors.
  Use permitted under Creative Commons License Attribution 4.0
  International (CC BY 4.0).}

\conference{LLM4SE 2026: Workshop on Large Language Models for Software Engineering,
  July 2026}

\title{EvolveTool-Bench: Diagnosing Tool Library Quality in Self-Evolving LLM Agents}

\author[1]{Alibek T. Kaliyev}[%
email=alibek.kaliyev@utexas.edu,
url=https://arise-ai.dev,
]
\address[1]{University of Texas at Austin, Austin, TX, USA}

\begin{abstract}
A growing body of work enables LLM agents to create their own tools at runtime---VOYAGER synthesizes code skills, EvoSkill evolves strategy-level skill folders, and ARISE generates executable Python tools with sandbox validation. Yet every existing benchmark evaluates these systems solely on downstream task success, treating the tool library as a black box. We introduce \textsc{EvolveTool-Bench}, the first diagnostic benchmark that evaluates the \emph{quality of the tool library itself}: redundancy, reuse, composition ability, regression stability, and safety. Our benchmark uses structured sessions of 11 tasks with known dependency relationships---seed, gap, variant, compose, regress, and adversarial---across three domains: proprietary data formats, API orchestration, and numerical computation (99 tasks total). We evaluate four systems---no-evolution, ARISE (code-level), EvoSkill (strategy-level), and one-shot creation---with two models (Claude Sonnet 4 and Claude Haiku 4.5). Our key findings: (1)~all six configurations achieve similar task completion (63--68\%), but EvolveTool Scores range from 0.516 to 0.612---library quality is the dominant differentiator invisible to task-completion-only evaluation; (2)~code-level evolution (ARISE, ETS=0.603) substantially outperforms strategy-level evolution (EvoSkill, ETS=0.520) and one-shot creation (ETS=0.519) in the first head-to-head comparison of these paradigms; (3)~iteration and sandbox validation are essential---one-shot creation produces 3 tools vs ARISE's 22, with negligible library health improvement; (4)~tool reuse scales with library quality, from 40.7\% (no evolution) to 74.1\% (ARISE).
\end{abstract}

\begin{keywords}
  LLM agents \sep
  tool synthesis \sep
  self-evolving agents \sep
  benchmark \sep
  tool library quality \sep
  software engineering
\end{keywords}

\maketitle

\section{Introduction}

LLM agents increasingly create their own tools at runtime. VOYAGER~\cite{wang2023voyager} synthesizes code skills in Minecraft, CREATOR~\cite{qian2023creator} uses tool creation as a reasoning strategy, LLMs as Tool Makers~\cite{cai2023latm} separates tool-making from tool-using, and EvoSkill~\cite{alzubi2026evoskill} evolves structured skill folders via failure analysis. ARISE (Adaptive Runtime Improvement through Self-Evolution)~\cite{kaliyev2026arise} takes a different approach: it sits between an agent and its tool library as framework-agnostic middleware, observing task failures, detecting capability gaps via LLM analysis, synthesizing executable Python tools, validating them through sandboxed execution and adversarial testing, and promoting passing tools to the agent's active library---all without human intervention. ARISE supports iterative refinement, A/B testing of tool variants, version-controlled rollback, and distributed deployment via S3/SQS.

These systems are evaluated exclusively on \emph{task success}: if the agent solves downstream tasks, the tool library is deemed good. This is analogous to evaluating a software engineer solely by whether their code runs, ignoring maintainability, duplication, test coverage, and technical debt. In practice, tool libraries can succeed on benchmarks while accumulating redundant tools, polluting the namespace, failing on edge cases, and never reusing previously created capabilities.

We present \textsc{EvolveTool-Bench}, a diagnostic benchmark designed to fill this evaluation gap. Rather than measuring only whether evolved tools help solve tasks, we measure \emph{how well agents build, maintain, and reuse tool libraries over time}. Our contributions:

\begin{enumerate}
\item \textbf{Library-level metrics.} Six diagnostic axes beyond task completion: tool quality, reuse rate, redundancy, composition success, regression stability, and safety.

\item \textbf{Controlled task sessions.} Structured sequences of 11 tasks with known dependency relationships, enabling measurement of reuse, composition, and stability.

\item \textbf{Three domains requiring tool execution.} Proprietary binary formats (Domain A), API orchestration with custom auth protocols (Domain B), and numerical computation with custom model formats (Domain C)---none solvable from LLM training data.
\end{enumerate}

\section{Related Work}

\paragraph{Tool Creation Benchmarks.} Tool-Genesis~\cite{xia2026toolgenesis} evaluates one-shot tool creation with surface compliance, schema-F1, and unit test pass rates---all per-tool metrics with no library-level evaluation. SkillsBench~\cite{skillsbench2026} evaluates whether curated skills improve task success but uses static skill sets, not evolved ones.

\paragraph{Self-Evolving Agent Frameworks.} VOYAGER~\cite{wang2023voyager} demonstrates lifelong skill acquisition in Minecraft, evaluated by exploration metrics (items discovered, distance traveled). EvoSkill~\cite{alzubi2026evoskill} evolves textual skill folders---structured prompts with trigger metadata and helper scripts---through failure-driven feedback with Pareto frontier selection, evaluated by downstream accuracy on QA benchmarks. ARISE~\cite{kaliyev2026arise} operates at the code level: it synthesizes executable Python functions with auto-generated test suites, validates them in isolated sandboxes (subprocess or Docker), runs adversarial testing to probe edge cases, and promotes only passing tools. ARISE supports incremental patching of existing tools, A/B testing between tool variants, and distributed deployment where stateless agents read tools from S3 while a background worker handles evolution. None of these systems evaluate tool library quality---all report only task success.

\paragraph{Tool Use Benchmarks.} ToolBench~\cite{qin2023toolbench}, AgentBench~\cite{liu2024agentbench}, and ToolComp~\cite{toolcomp2025} evaluate agents using pre-defined tool sets. Vending-Bench~\cite{backlund2025vendingbench} tests long-horizon coherence with fixed tools. None evaluate agents that create their own tools.

\paragraph{Gap.} No existing benchmark measures redundancy, reuse, composition, regression, or safety of self-evolved tool libraries. Our work fills this gap.

\section{Benchmark Design}

\subsection{Task Session Structure}

Each session contains 11 tasks in six categories with known dependency relationships (Table~\ref{tab:task_types}). Variant and regress tasks are designed so that the evaluator \emph{knows} which prior tool should be reused, enabling measurement of reuse rate independent of task success.

\begin{table}[t]
\centering
\small
\begin{tabular}{llp{4.5cm}}
\toprule
\textbf{Type} & \textbf{N} & \textbf{Diagnoses} \\
\midrule
Seed & 3 & Can agent use provided tools? \\
Gap & 2 & Does agent create a working tool? \\
Variant & 2 & Reuse or create duplicate? \\
Compose & 1 & Can agent chain own tools? \\
Regress & 1 & Does the library stay stable? \\
Adversarial & 2 & Can agent handle edge cases? \\
\bottomrule
\end{tabular}
\caption{Task types in each diagnostic session.}
\label{tab:task_types}
\end{table}

\subsection{Domains}

\paragraph{Domain A: Data Transformation.} Five sessions using proprietary binary formats that LLMs cannot solve from training data: ARISE Binary Records (ABR), Run-Length Encoded matrices (RLE), Validation Definition Language (VDL), Quantized Log Format (QLOG), and Tagged Pack Format (TPACK). Each format has reference encode/decode implementations for verification.

\paragraph{Domain B: API Orchestration.} Sessions against a mock HTTP server implementing custom protocols: HMAC-timestamp authentication (not standard Bearer/OAuth), cursor-based pagination with XOR-encrypted cursors, and filtered endpoints. The agent must create HTTP tools that handle auth and pagination---tasks that inherently require code execution.

\paragraph{Domain C: Numerical Computation.} Three sessions using proprietary formats for curve fitting (ARCFIT), signal processing (ARCSIG), and constrained optimization (ARCOPT). Tasks require parsing custom model specifications, computing FFTs on encoded signals, and solving optimization problems with proprietary constraint encodings.

\subsection{Evaluation Axes}

\paragraph{Axis 1: Task Completion.} Binary pass/fail via deterministic verification. Reported overall and by task type.

\paragraph{Axis 2: Tool Quality Score (TQS).} Each synthesized tool is evaluated on four dimensions, each scored 0--1:

\begin{itemize}
\item \textbf{Correctness}: Executed on $N$ hidden unit tests (not shown to agent). Score = proportion passing.
\item \textbf{Robustness}: Executed on $M$ adversarial inputs (empty strings, null values, type errors). Score = proportion handled without crashing.
\item \textbf{Generality}: Executed on held-out inputs from the same distribution. Score = proportion producing valid output.
\item \textbf{Code Quality}: Automated static checks---docstring, type hints, exception handling, no magic numbers, under 100 lines, valid syntax.
\end{itemize}

\begin{equation}
\text{TQS}(t) = \frac{1}{4}\bigl(C(t) + R(t) + G(t) + Q(t)\bigr)
\label{eq:tqs}
\end{equation}
where $C$ is correctness, $R$ robustness, $G$ generality, and $Q$ code quality, each in $[0,1]$. All evaluation runs in isolated subprocesses with 10-second timeouts.

\paragraph{Axis 3: Library Diagnostics.} Six metrics measured at session end:
\begin{align}
\rho_r &= \frac{|\{t \in \mathcal{T}_{v,r} : \text{reused}\}|}{|\mathcal{T}_{v,r}|} & \text{(Reuse Rate)} \label{eq:reuse} \\[2pt]
\rho_d &= \frac{|\{(t_i,t_j) : \text{output\_equiv}\}|}{\binom{|\mathcal{L}|}{2}} & \text{(Redundancy)} \label{eq:redund} \\[2pt]
\pi &= \frac{|\{t \in \mathcal{L} : \text{TQS}(t) \geq 0.5\}|}{|\mathcal{L}|} & \text{(Precision)} \label{eq:prec} \\[2pt]
\eta &= \frac{|\mathcal{L}_{\text{used}}|}{|\mathcal{L}|} & \text{(Efficiency)} \label{eq:eff} \\[2pt]
\gamma &= \frac{|\{t \in \mathcal{T}_c : \text{passed}\}|}{|\mathcal{T}_c|} & \text{(Composition)} \label{eq:comp} \\[2pt]
\delta &= \frac{|\{t \in \mathcal{T}_{reg} : \text{failed}\}|}{|\mathcal{T}_{reg}|} & \text{(Regression)} \label{eq:reg}
\end{align}
where $\mathcal{T}_{v,r}$ are variant/regress tasks, $\mathcal{L}$ is the tool library, $\mathcal{T}_c$ are compose tasks, and $\mathcal{T}_{reg}$ are regress tasks.

\paragraph{Library Health (LH).} Aggregates the six diagnostics:
\begin{equation}
\text{LH} = \frac{1}{6}\bigl(\rho_r + (1{-}\rho_d) + \pi + \eta + \gamma + (1{-}\delta)\bigr)
\label{eq:lh}
\end{equation}

\paragraph{Composite Score.} The EvolveTool Score (ETS) combines all axes:
\begin{equation}
\text{ETS} = \alpha_1 \cdot \text{TC} + \alpha_2 \cdot \overline{\text{TQS}} + \alpha_3 \cdot (1{-}\text{RC}) + \alpha_4 \cdot \text{LH} + \alpha_5 \cdot \text{SS}
\label{eq:ets}
\end{equation}
with $\alpha = (0.25, 0.20, 0.10, 0.30, 0.15)$. TC is task completion, $\overline{\text{TQS}}$ is mean tool quality, RC is refinement cost, LH is library health (Eq.~\ref{eq:lh}), and SS is safety score. Library health receives the highest weight as our primary contribution.

\section{Experimental Setup}

\paragraph{Systems.} Four systems testing different evolution strategies:
\begin{enumerate}
\item \textbf{No-Evolution}: Seed tools only, no tool creation ability.
\item \textbf{ARISE}: Full iterative evolution with sandbox validation, adversarial testing, and LLM-as-judge reward. Failure threshold 1.
\item \textbf{EvoSkill}: Strategy-level evolution---generates text strategies (prompts), not executable code.
\item \textbf{One-Shot}: Creates tools on first failure with no sandbox testing, no adversarial validation, no refinement.
\end{enumerate}

\paragraph{Models.} Claude Sonnet 4 and Claude Haiku 4.5, both via Amazon Bedrock.

\paragraph{Benchmark.} Domain A: 5 sessions $\times$ 11 tasks = 55. Domain B: 1 session $\times$ 11 = 11. Domain C: 3 sessions $\times$ 11 = 33. Total: \textbf{99 tasks} per system-model configuration.

\section{Results}

\begin{table*}[t]
\centering
\small
\begin{tabular}{llcccccc}
\toprule
\textbf{System} & \textbf{Model} & \textbf{ETS}$\uparrow$ & \textbf{TC (\%)}$\uparrow$ & \textbf{Tools} & \textbf{TQS}$\uparrow$ & \textbf{Reuse (\%)}$\uparrow$ & \textbf{LH (\%)}$\uparrow$ \\
\midrule
No-Evolution & Sonnet & .518 & 66.7 & 0 & --- & 48.1 & 41.4 \\
EvoSkill & Sonnet & .520 & \textbf{68.2} & 0 & --- & 48.1 & 41.4 \\
One-Shot & Sonnet & .519 & 67.7 & 3 & .04 & 48.1 & 38.3 \\
\textbf{ARISE} & \textbf{Sonnet} & .603 & 63.6 & \textbf{22} & .34 & \textbf{70.4} & \textbf{51.7} \\
\midrule
No-Evolution & Haiku & .516 & 66.7 & 0 & --- & 51.9 & 40.1 \\
\textbf{ARISE} & \textbf{Haiku} & \textbf{.612} & 65.2 & 13 & \textbf{.35} & 66.7 & 51.2 \\
\bottomrule
\end{tabular}
\caption{Full benchmark results (99 tasks per configuration, 3 domains). ARISE achieves highest ETS on both models despite slightly lower task completion. ARISE+Haiku creates fewer tools (13 vs 22) but achieves the highest overall ETS (0.612).}
\label{tab:results_full}
\end{table*}

\subsection{Cross-System Analysis}

Table~\ref{tab:results_full} presents results across all six configurations. The central finding: \textbf{task completion is similar across systems (63--68\%), but library quality diverges dramatically.} ARISE's ETS (0.603--0.612) is 16--18\% higher than all other systems ($\sim$0.52), driven entirely by library-level metrics.

\paragraph{Code vs.\ Strategy Evolution.} ARISE (code-level) and EvoSkill (strategy-level) represent fundamentally different approaches. EvoSkill's text strategies modestly improve task completion (68.2\% vs 66.7\%) but create zero executable tools and achieve identical library health to no-evolution (41.4\%). ARISE creates 22 executable tools with 70.4\% reuse. A hybrid approach combining strategy guidance with code synthesis may be optimal.

\paragraph{Iteration vs.\ One-Shot.} One-Shot creates 3 tools without validation. ARISE creates 22 with sandbox testing and adversarial validation. Despite similar task completion, One-Shot's library health (38.3\%) is \emph{below} no-evolution (41.4\%). Unvalidated tools are unreliable.

\paragraph{Example: Evolved Tool.} When ARISE's agent failed to authenticate with the mock API's HMAC-timestamp scheme, it synthesized an \texttt{hmac\_sha256} tool:

\begin{lstlisting}[language=Python, basicstyle=\small\ttfamily]
def hmac_sha256(secret: str,
                message: str) -> str:
    """Generate HMAC-SHA256 hash."""
    import hmac, hashlib
    try:
        h = hmac.new(
            secret.encode('utf-8'),
            message.encode('utf-8'),
            hashlib.sha256)
        return h.hexdigest()
    except Exception as e:
        return f"Error: {str(e)}"
\end{lstlisting}

Generated, tested (5/5 sandbox tests passed), and promoted in $<$60 seconds.

\subsection{Key Findings}

\paragraph{Finding 1: Task completion alone is misleading.} All six configurations achieve similar task completion (63--68\%), yet ETS scores range from 0.516 to 0.612. ARISE achieves the \emph{lowest} task completion (63.6\%) but the \emph{highest} ETS. \textbf{A benchmark reporting only TC would rank ARISE last---the opposite of the library quality picture.}

\paragraph{Finding 2: Code evolution outperforms strategy evolution.} ARISE (ETS=0.603) outperforms EvoSkill (ETS=0.520) in the first head-to-head comparison. EvoSkill creates zero executable tools. This is the first empirical evidence that code-level evolution produces better tool libraries than strategy-level evolution.

\paragraph{Finding 3: Synthesis without validation is insufficient.} One-Shot creates 3 tools vs ARISE's 22. Its library health (38.3\%) is below no-evolution (41.4\%). Without sandbox testing and adversarial validation, synthesized tools are unreliable.

\paragraph{Finding 4: Reward design is the critical bottleneck.} In preliminary experiments with heuristic rewards, ARISE created zero tools. The agent produces helpful refusal responses that score 1.0 under heuristic evaluation. Only LLM-as-judge rewards correctly identify failures. \textbf{The synthesis pipeline is not the bottleneck---the reward signal is.}

\paragraph{Finding 5: Tool quality is the frontier.} ARISE's tools average TQS=0.34--0.35, with correctness being the weakest dimension. Many synthesized tools parse proprietary formats with off-by-one errors or incomplete edge case handling. The benchmark reveals this via hidden tests the agent never sees.

\paragraph{Finding 6: Cheaper models match library quality.} ARISE+Haiku achieves the highest ETS (0.612) with 13 tools, outperforming ARISE+Sonnet (0.603, 22 tools). Tool evolution does not require the most capable model.

\section{Discussion}

\paragraph{The Benchmark Contribution.} The primary contribution is the benchmark itself, not the performance of any particular system. The controlled session structure (seed$\to$gap$\to$variant$\to$compose$\to$regress$\to$adversarial) with known dependencies is reusable across any domain.

\paragraph{Honest Assessment of ARISE.} Our results expose both strengths and limitations. Strengths: ARISE is the only system that builds meaningful, reusable tool libraries (22 tools, 70.4\% reuse). Limitations: (1)~evolution overhead reduces task completion below no-evolution; (2)~tool quality averages only TQS=0.34; (3)~complex binary formats often exceed the synthesis LLM's ability to generate correct parsers.

\paragraph{ETS Design Choices.} Systems creating zero tools receive TQS=0 in ETS. This is intentional: \textsc{EvolveTool-Bench} evaluates tool library quality, and systems that cannot create tools should score lower. ETS is not a general agent ranking; for that, TC should be used alongside ETS.

\paragraph{Limitations.} (1)~Our EvoSkill baseline approximates the original system; running the actual EvoSkill codebase would strengthen the comparison. (2)~Two models from one provider (Anthropic via Bedrock); open-source models would broaden generalizability. (3)~Tool quality uses automated checks; human evaluation would capture nuances our metrics miss. (4)~No statistical significance testing across seeds.

\section{Conclusion}

We introduce \textsc{EvolveTool-Bench}, the first benchmark to evaluate self-evolved tool libraries as first-class artifacts. Across 99 tasks in three domains, we find that task completion---the standard evaluation metric---fails to differentiate systems with fundamentally different library quality. Code-level evolution (ARISE) builds libraries with 70.4\% reuse and 51.7\% health on Sonnet, and achieves even higher ETS (0.612) with the cheaper Haiku model, outperforming strategy-level evolution (EvoSkill) and one-shot creation, but at the cost of lower task completion and imperfect tool quality. The benchmark reveals that reward design, not synthesis capability, is currently the critical bottleneck for self-evolving systems.

We release the benchmark, evaluation harness, proprietary format specifications, and all experimental data.\footnote{Code: \url{https://github.com/abekek/evolvetool-bench}}

\section*{Declaration on Generative AI}
During the preparation of this work, the author used Claude (Anthropic) for code generation assistance in implementing the benchmark harness, baseline systems, and experimental infrastructure. The author reviewed and edited all content and takes full responsibility for the publication's content. All experimental results were produced by automated benchmark runs without manual intervention.

\begin{acknowledgments}
The author thanks the Strands Agents and ARISE open-source communities for infrastructure support.
\end{acknowledgments}

\bibliography{references}

\end{document}
